\begin{abstract}
We propose News Event Decay (NED), a neural architecture for
high-frequency cryptocurrency volatility forecasting that represents
news through article-dependent magnitudes and decay rates. Pretrained
CryptoBERT embeddings determine these parameters, and their decaying
signals form a bounded correction to a neural time-series variance
forecast. Across six expanding-window test periods from 2023 to 2025,
NED reduces Bitcoin QLIKE by 7.38\% at the one-hour horizon and
10.83\% at four hours, with log-MSE reductions of 0.99\% and 3.11\%
($p<0.01$). Ethereum QLIKE gains are smaller and its log-MSE differences
are insignificant. Bitcoin improvements are strongest in the tails of
the subsequent standardized-return distribution. In BTC four-hour
models, the article half-life has a median of 42.7 minutes, and negative-classified
articles receive larger average signal magnitudes than
positive-classified articles. Macro news receives the largest mean
magnitudes and longest median half-lives.
Matched experiments favor learned decay over constant aggregation,
although the constant kernel also improves forecasts. Sentiment
with learned decay nearly matches the full embedding, while PCA
compression improves on it in this run.
\end{abstract}

\noindent\textbf{Keywords:} Realized volatility; Volatility forecasting;
News-driven volatility; Cryptocurrency; Hawkes processes; Text embeddings.

\noindent\textbf{JEL Codes:} C45, C53, C58, G17.
